%% ============================================================
%% AULA 8 — Redes que leem: sequências, recorrência e memória
%% GBC073 — Inteligência Computacional (FACOM/UFU)
%%
%% Esqueleto com esboço de conteúdo — a desenvolver.
%% Ementa (ficha SEI 5116656): item 2 — arquiteturas;
%% deep learning: introdução.
%%
%% Compilar:  latexmk -xelatex aula08.tex   (duas passadas!)
%% ============================================================
\documentclass[aspectratio=169, 11pt]{beamer}

\makeatletter
\def\input@path{{./}{../}}
\makeatother

\usetheme{InteligenciaComputacional}   % ANTES do babel: fontspec vem primeiro
\usepackage{babel}
\babelprovide[import, main]{brazilian}

\title[Aula 8 — Recorrência]{Aula 8: Redes que leem — sequências, recorrência e memória}
\subtitle[GBC073 — Inteligência Computacional]{RNN \textbullet\ gradiente que some \textbullet\ LSTM e GRU}
\author[Prof.\ Marcelo Albertini]{Prof.\ Marcelo Keese Albertini}
\institute{GBC073 — Inteligência Computacional \textbullet\ Universidade Federal de Uberlândia}
\date{2026}

\begin{document}

% ------------------------------------------------------------ capa
\begin{frame}[plain, noframenumbering]
  \titlepage
\end{frame}

% ------------------------------------------------------------ roteiro
\begin{frame}{Roteiro da aula}
  \begin{itemize}\setlength{\itemsep}{0.5ex}
    \item<1-> \textbf{Sequências não cabem em MLPs}
    \item<2-> \textbf{RNN e o gradiente que some (ou explode)}
    \item<3-> \textbf{LSTM e GRU: memória com portões}
    \item<4-> \textbf{Aplicações: linguagem e séries temporais}
  \end{itemize}
\end{frame}

% ------------------------------------------------------------
\section{Sequências não cabem em MLPs}

\begin{frame}{Sequências não cabem em MLPs}
  \begin{itemize}\setlength{\itemsep}{0.4ex}
    \item Texto, áudio, séries temporais: comprimento variável, ordem importa.
    \item Uma MLP exige tamanho fixo de entrada — e ignora a ordem.
    \item A recorrência: o estado oculto carrega o passado de um passo ao
          próximo.
    \item A mesma célula é reaplicada a cada passo — pesos compartilhados no
          tempo.
  \end{itemize}
\end{frame}

% ------------------------------------------------------------
\section{RNN e o gradiente que some}

\begin{frame}{RNN e o gradiente que some (ou explode)}
  \begin{itemize}\setlength{\itemsep}{0.4ex}
    \item Desenrolar a RNN no tempo (\emph{unrolling}): uma rede profunda por
          passo.
    \item Retropropagação através do tempo (BPTT).
    \item Produtos repetidos de jacobianos: o gradiente some (ou explode)
          exponencialmente.
    \item O gradiente que some apaga dependências de longo alcance.
  \end{itemize}
\end{frame}

% ------------------------------------------------------------
\section{LSTM e GRU}

\begin{frame}{LSTM e GRU: memória com portões}
  \begin{itemize}\setlength{\itemsep}{0.4ex}
    \item LSTM: célula de memória $c$ com portões de esquecer, escrever e ler.
    \item O portão de esquecer abre um caminho de gradiente que não some.
    \item GRU: versão enxuta — dois portões, sem célula separada.
    \item O que as duas fazem de verdade: escolher o que lembrar e o que
          descartar.
  \end{itemize}
\end{frame}

% ------------------------------------------------------------
\section{Aplicações}

\begin{frame}{Aplicações: linguagem e séries temporais}
  \begin{itemize}\setlength{\itemsep}{0.4ex}
    \item Modelagem de linguagem: prever o próximo símbolo.
    \item Tradução, classificação de sentimento, legendagem.
    \item Séries temporais: previsão, anomalias, controle.
    \item Na prática moderna, muitas tarefas migraram para atenção (aula 9) —
          mas a recorrência continua viva em nichos.
  \end{itemize}
\end{frame}

\end{document}
